\chapter{Introduction}
\label{Chapter 1}

\section{Introduction}
Retail lending in India has exploded over the past ten years. Personal loans. Two-wheeler loans. Consumer-durable EMIs, gold loans, credit-card revolving balances, salary advances, and the whole buy-now-pay-later wave that came out of the pandemic. Every kind of small-ticket loan is now being pushed by public-sector banks, private banks, small finance banks, NBFCs and a long list of digital-only fintech lenders that didn't exist five years ago. With that volume comes the other side of the ledger: people who miss payments. A small but real share of every cohort fails to pay an EMI on time, and the lender has to chase them. Partly because the RBI expects it. Partly because letting an account drift turns expensive very fast.

The chasing itself is what the industry calls ``collections''. Once an account goes even a day past due (DPD), it enters a standard escalation pipeline, with the tone of follow-up getting harder the longer the dues stay unpaid. Table~\ref{tab:dpd-buckets} shows the bucketing that Indian retail lenders broadly agree on.

\begin{table}[h]
\centering
\caption{Collection buckets typical to Indian retail-lending portfolios.}
\label{tab:dpd-buckets}
\begin{tabular}{lll}
\toprule
\textbf{Bucket} & \textbf{DPD range} & \textbf{Typical activity} \\
\midrule
Bucket 0 & 1--30 days   & Soft reminders, automated calls, light SMS \\
Bucket 1 & 31--60 days  & Tele-calling, WhatsApp, agent calls \\
Bucket 2 & 61--90 days  & Field visits, legal notice preparation \\
Bucket 3 & 91--120 days & Pre-write-off, recovery escalation \\
NPA      & 120+ days    & Settlement, legal action \\
\bottomrule
\end{tabular}
\end{table}

Recovery gets dramatically harder as the bucket deepens. An account cleared inside Bucket~0 is cheap to recover, takes almost no negotiation, and never forces the lender to set aside a provision. An account that slips down into Bucket~3 is the opposite story -- expensive to chase, almost always settled at a haircut, and the customer relationship rarely survives. So every lender wants to clean up Bucket~0 aggressively. Easier said than done. A lender holding a hundred thousand active retail loans can comfortably have eight to twelve thousand accounts sitting in Bucket~0 on any given day, and each one needs a contact attempt inside a tight window. The usual answer is outsourcing -- a recovery agency takes the chase on, keeps a cut of whatever it brings in, and hands the rest back to the lender.

The work in this report came out of an internship at Ignosis, a fintech infrastructure company best known in the Account Aggregator (AA) ecosystem. Over the last year and a half, the firm has layered a collections product on top of its existing AA stack, reusing the same multi-tenant infrastructure and folding AA-pulled bank-statement data into how cases get prioritised. I joined as a comprehensive-project intern on 1 January 2026 and stayed full-time until 31 May 2026. My work sat across two repositories: the orchestrator backend (plus its dashboards) on one side, and a separate automation toolkit holding the WhatsApp service, the desktop client and the Android SMS sender on the other. The project's title is \textbf{``AI-Based Collections Platform''}. Architecturally, it is a multi-tenant web application: a case management dashboard for agents, a bulk campaign module for managers, an autonomous AI voice caller, outreach across WhatsApp and SMS, and reporting on contact rate, agent productivity, call dispositions and the recovery number that everyone watches.

\section{Motivation}
The project lives at the junction of a real industry pain point and a technology that has only recently become practical. The industry side is straightforward. Every percentage point of Bucket~0 recovery moves the lender's Gross Non-Performing Asset (GNPA) ratio and translates into less money set aside for provisioning. Pushing the early-bucket rate from the usual 60--70 percent into the mid-80s, which is what the Ignosis customer base now expects, has a direct rupee value for any lender of meaningful size.

The technology side is where things have actually changed. Three shifts over the past couple of years made a project like this even feasible, none of which were really in place in early 2023. Conversational voice models have crossed a real usability threshold on Indian English and on most of the regional languages, to the point that an AI agent can hold an outbound conversation about a payment due date and the customer on the other end mostly cannot tell it isn't a person; the firm's internal voice AI platform runs these agents at latencies that surprised me the first time I watched a live call. Account Aggregator data is the second shift -- a lender can now pull bank statements with consent through the AA framework, which lifts the opening of the conversation from a generic ``please pay'' to a much more specific ``we noticed your salary credit came in yesterday, what date works for you''. The third shift is operational: it is now feasible to run a software-controlled pool of consumer SIM cards, sending WhatsApp from real phone numbers and SMS from real Android phones, which sidesteps the deliverability ceiling that business APIs cannot escape because of their pricing and approval flows.

Putting these three capabilities behind one dashboard, with case management and bulk campaigns wrapped around them, was the engineering goal of the internship.

\section{Objective}
Going into the internship I had six concrete objectives. All six were either fully met or substantially met by the end of May.

The first was to build a centralised case management interface, accessible to agents in a browser, that would show every allocated case together with its payment history and prior interaction notes, and let an agent act on the case directly without context-switching. The second was a bulk campaign module: the campaign manager should be able to wrap a batch of cases -- typically a CSV that lands at the start of the monthly cycle -- into one schedulable unit. The third was to wire AI-powered voice calling into that campaign module, so customer outreach could continue after the floor of agents had signed off for the day. The fourth was a multi-channel follow-up layer across WhatsApp and SMS that ran on real consumer SIMs instead of business APIs, with anti-block heuristics and automatic device rotation. The fifth was an analytical pass on around forty thousand customer records, looking for behavioural patterns that should inform prioritisation, and folding those findings back into the workflow templates that ship by default. The sixth was procedural rather than functional: throughout the internship, hold to the engineering practices of the host company -- code review, unit and integration tests, conventional commits, ticketed work items and zero-downtime deployments.

\section{Problem Statement}
Without a centralised platform, a typical Indian retail lender's collections function runs into a handful of structural problems that reinforce each other.

It starts with case allocation. The lender drops a CSV or Excel of delinquent accounts on the agency at the start of every cycle, and the agency manager then distributes those cases by hand across whichever agents are on the floor that morning. The mapping is opaque to the lender, hard to audit later, and not really reproducible. The outreach itself is single-channel by default. Agents reach for the desk phone first. If the customer does not pick up, sometimes a SMS goes out and often it does not, and there is no orchestrated follow-up across phone, WhatsApp and SMS, so a real share of customers never gets reached on the channel they would have actually responded to. A surprising slice of every agent's day is also spent on what is essentially overhead: finding the next case in a spreadsheet, dialling the number manually, skimming the customer's history before the call. The actual conversation is the smaller half of the hour.

On top of that, there are two structural blind spots. The first is the absence of after-hours coverage. An agent's day ends at 6 or 7 PM and after that the lender effectively goes dark, even though many salaried customers are easiest to reach in the late evening when they are home from work. The second is measurement. Recovery rates surface as a single summary number at month-end, with no real way for lenders to look at bucket-level conversion, channel-level effectiveness or agent-level productivity. And underneath both is a deliverability ceiling. WhatsApp and SMS gateways via business APIs throttle hard, charge per message, and require template approvals that drag on for days, while SIM-based outreach from real consumer numbers reaches customers much more reliably. No off-the-shelf product gives a collections team that second option.

In compact form the problem statement is:

\textit{``Design and contribute to a centralised, multi-tenant collections platform that improves contact rate, agent throughput and after-hours coverage by integrating AI voice calling and multi-channel outreach across WhatsApp and SMS into a single case management workflow, while preserving end-to-end auditability and supporting bulk campaign execution at the scale of a typical monthly allocation.''}

\section{Approach}
The way we worked was iterative and code-first. Instead of writing a long design document up front and then implementing it, the team follows what the engineering lead calls plan-first-but-small: every piece of work starts as a short markdown plan attached to a ticket, the plan is reviewed by the lead (Mr.~Rahi Shah in my case), a branch is cut, and something working is on a dev environment inside a few days. Production deployment happens only after code review, integration testing on a staging tenant and a verification pass by QA. Every artefact -- the plan, the branch, the merge request, the QA notes -- stays in the repo so that the trail of decisions can still be reconstructed months later. This is the dominant cultural fact about how the work in this report got done, and it explains why my contributions are spread across several branches that build on each other rather than sitting in one big feature branch.

\section{Scope of the Project}
The scope of the internship covers:
\begin{itemize}
\item Backend contributions to the multi-tenant orchestrator monolith: data entities, REST controllers, services, scheduled jobs and schema migrations.
\item Frontend contributions to the agent dashboard for case management, campaign execution and analytics.
\item A separate automation toolkit containing the WhatsApp service, the desktop application, the Android SMS sender and an internal REST API server.
\item A one-off analysis on customer-level transaction and repayment data.
\end{itemize}

The scope deliberately leaves out the AI voice agent itself (which is the firm's internal voice AI platform), the telephony gateway, the cloud infrastructure provisioning, the lender's internal LMS, and the regulatory compliance review.

\section{Organization of the Rest of the Report}

The rest of the report is organised as follows.

\textbf{Chapter 2 -- Literature Review.} A look at the relevant prior art: collection-management products in India and globally, the conversational voice agent space, the WhatsApp multi-device protocol, and the academic work on anti-spam and rate-limiting in messaging platforms. The chapter ends by translating the problem statement of Chapter~1 into a set of concrete engineering requirements.

\textbf{Chapter 3 -- Proposed Methodology and Implementation.} A detailed walk through the development methodology used at Ignosis and through each of the modules I worked on -- the case management subsystem, the workflow engine, the campaign module, the AI voice calling integration, the WhatsApp automation, the desktop application and the Android SMS sender.

\textbf{Chapter 4 -- Software Design.} A design-level view: goals, component decomposition, data model, API surface, the principal sequence flows, the deployment view and the security design.

\textbf{Chapter 5 -- Results and Discussion.} The quantitative and qualitative results from running the platform in production: the early-bucket collection rate, the percentage of cases recovered inside the first five days, the contact-rate uplift from multi-channel outreach, the throughput of the AI voice agent, and the WhatsApp ban rates observed in practice.

\textbf{Chapter 6 -- Conclusions and Future Scope.} What was achieved, what I would do differently in hindsight, and what the team has lined up for the next two quarters.

The References and an Appendix with platform reference material close out the document.
